\documentclass[11pt,letterpaper]{article}
% This source is intentionally restricted to pdfLaTeX.
\usepackage[T1]{fontenc}
\usepackage[utf8]{inputenc}
\usepackage{textcomp}
\usepackage{lmodern}
\usepackage[margin=0.82in,top=0.94in,bottom=0.91in,headheight=18pt]{geometry}
\usepackage{amsmath,amssymb}
\usepackage[table]{xcolor}
\definecolor{DeepNavy}{HTML}{14283F}
\definecolor{Teal}{HTML}{0A747C}
\definecolor{Gold}{HTML}{C78F36}
\definecolor{Slate}{HTML}{536477}
\definecolor{Pale}{HTML}{F4F7FA}
\definecolor{CodeEdge}{HTML}{D4DFE9}
\usepackage{microtype}
\usepackage{parskip}
\setlength{\parskip}{0.47em}
\usepackage{titlesec}
\titleformat{\section}{\normalfont\Large\bfseries\color{DeepNavy}}{\thesection}{0.62em}{}
\titlespacing*{\section}{0pt}{2.0ex plus .7ex minus .2ex}{0.75ex}
\titleformat{\subsection}{\normalfont\large\bfseries\color{Teal}}{\thesubsection}{0.6em}{}
\usepackage{longtable,booktabs,array,calc}
\renewcommand{\arraystretch}{1.22}
\AtBeginEnvironment{longtable}{\small}
\usepackage{etoolbox}
\makeatletter
\patchcmd\longtable{\par}{\if@noskipsec\mbox{}\fi\par}{}{}
\makeatother
\usepackage{listings}
\IfFileExists{upquote.sty}{\usepackage{upquote}}{}
\lstdefinestyle{suitecode}{
  basicstyle=\ttfamily\footnotesize,
  breaklines=true,breakatwhitespace=false,
  columns=fullflexible,keepspaces=true,showstringspaces=false,
  frame=single,rulecolor=\color{CodeEdge},
  backgroundcolor=\color{Pale},
  keywordstyle=\color{Teal}\bfseries,
  commentstyle=\color{Slate}\itshape,
  stringstyle=\color{DeepNavy},
  xleftmargin=5pt,xrightmargin=5pt,
  framexleftmargin=5pt,framexrightmargin=5pt,
  aboveskip=0.85em,belowskip=0.95em,
  tabsize=2
}
\lstset{style=suitecode}
\usepackage{enumitem}
\setlist[itemize]{leftmargin=1.6em,itemsep=2pt,topsep=3pt}
\setlist[enumerate]{leftmargin=1.8em,itemsep=2pt,topsep=3pt}
\providecommand{\tightlist}{\setlength{\itemsep}{0pt}\setlength{\parskip}{0pt}}
\newcommand{\passthrough}[1]{#1}
\usepackage{xurl}
\usepackage[colorlinks=true,linkcolor=Teal,urlcolor=Teal,citecolor=Teal,pdfborder={0 0 0}]{hyperref}
\urlstyle{same}
\usepackage{fancyhdr,lastpage}
\pagestyle{fancy}
\fancyhf{}
\fancyhead[L]{\footnotesize\sffamily\color{Slate} ECMA-262:2025 \textbullet{} Engineering Documentation}
\fancyhead[R]{\footnotesize\sffamily\color{Slate} 16th Edition}
\fancyfoot[L]{\footnotesize\sffamily\color{Slate} Independent technical companion}
\fancyfoot[R]{\footnotesize\sffamily\color{Slate}\thepage\ / \pageref{LastPage}}
\renewcommand{\headrulewidth}{0.35pt}
\renewcommand{\footrulewidth}{0pt}
\setlength{\emergencystretch}{3em}
\setcounter{tocdepth}{2}
\setcounter{secnumdepth}{2}
\begin{document}
\begin{titlepage}
\thispagestyle{empty}
\vspace*{2.8cm}
{\color{Gold}\rule{18mm}{3pt}}\par\vspace{18mm}
{\sffamily\Large\bfseries\color{Teal} SOFTWARE ENGINEERING DOCUMENTATION SUITE}\par
\vspace{8mm}
{\sffamily\huge\bfseries\color{DeepNavy} ECMAScript 2025 }\par
\vspace{4mm}
{\sffamily\fontsize{24}{29}\selectfont\bfseries\color{DeepNavy} Runtime and Execution Model }\par
\vspace{10mm}
{\large\color{Slate} Realms, jobs, agents and the memory model }\par
\vspace{16mm}
{\color{Teal}\rule{\textwidth}{0.8pt}}\par
\vfill
{\sffamily\color{Slate}\large ECMA-262:2025 \textbullet{} 16th Edition \textbullet{} June 2025}\par
\vspace{3mm}
{\sffamily\color{Slate} October 2026 \textbullet{} Document 4 of 8}\par
\vspace{3mm}
{\small\color{Slate} Original engineering companion. The uploaded ECMA-262:2025 specification remains the authoritative source.}\par
\end{titlepage}
\pagenumbering{arabic}
\tableofcontents
\clearpage
\section{Runtime model and vocabulary}\label{runtime-model-and-vocabulary}

ECMA-262:2025 does not specify one particular virtual machine implementation, optimizer, garbage collector, or browser event loop. It specifies observable language semantics using execution contexts, environments, realms, jobs, agents, and a shared-memory model. This technical reference maps those abstractions to architectural reasoning without asserting that all engines implement the internals literally.

\textbf{{[}SPEC{]}} = ECMA-262-defined; \textbf{{[}HOST{]}} = host integration; \textbf{{[}MODEL{]}} = explanatory mental model or engineering guidance. Principal sources: clauses 5--10, 15--16, 25--29 and Annex D.

\section{Execution lifecycle and conceptual layers}\label{execution-lifecycle-and-conceptual-layers}

\begin{longtable}[]{@{}
  >{\raggedright\arraybackslash}p{(\columnwidth - 6\tabcolsep) * \real{0.2500}}
  >{\raggedright\arraybackslash}p{(\columnwidth - 6\tabcolsep) * \real{0.2500}}
  >{\raggedright\arraybackslash}p{(\columnwidth - 6\tabcolsep) * \real{0.2500}}
  >{\raggedright\arraybackslash}p{(\columnwidth - 6\tabcolsep) * \real{0.2500}}@{}}
\toprule\noalign{}
\begin{minipage}[b]{\linewidth}\raggedright
Layer
\end{minipage} & \begin{minipage}[b]{\linewidth}\raggedright
Standard artifact
\end{minipage} & \begin{minipage}[b]{\linewidth}\raggedright
Responsibility
\end{minipage} & \begin{minipage}[b]{\linewidth}\raggedright
Not specified here
\end{minipage} \\
\midrule\noalign{}
\endhead
\bottomrule\noalign{}
\endlastfoot
Source input & Source Text / grammar (11--12) & recognizes tokens, syntax, early errors & file discovery and editor behavior \\
Static semantics & syntax-directed operations (8) & binding lists, contains, early errors & engine AST representation \\
Script/module representation & Script and Module Records (16) & instantiate/link/evaluate language code & host filesystem/URL resolution \\
Context & Execution Context (9.4) & running function/script/module evaluation & OS thread allocation \\
Environment & Environment Records (9.1) & mapping identifier bindings & physical stack layout \\
Realm & Realm Records (9.3) & associated intrinsics and global state & iframe creation details \\
Jobs & Job abstract closures (9.5) & serialized execution work items & browser task loop schedule \\
Agent & Agent model (9.6--9.8) & independent execution participation & web workers API \\
\end{longtable}

\section{Environment records and identifier resolution}\label{environment-records-and-identifier-resolution}

\textbf{{[}SPEC{]}} Declarative, object, function, global, and module environment records support different binding operations. \passthrough{\lstinline!ResolveBinding!}/\passthrough{\lstinline!GetIdentifierReference!} follow environment chains. Lexical declarations create bindings before initialization; uninitialized access results in a \passthrough{\lstinline!ReferenceError!}. \passthrough{\lstinline!var!} and lexical declarations have different instantiation procedures.

\begin{lstlisting}
function createAccumulator(seed) {
  let total = seed;
  return step => (total += step);  // captures environment binding
}
const add = createAccumulator(10);
console.assert(add(2) === 12);
console.assert(add(3) === 15);
\end{lstlisting}

\textbf{{[}MODEL{]}} Draw a closure as a function object associated with an outer lexical environment, not as a frozen copy of current variable values. In debugger traces distinguish creation time, initialization time, and evaluation time.

\section{Realms and intrinsics}\label{realms-and-intrinsics}

\textbf{{[}SPEC{]}} A realm encapsulates a set of intrinsic objects and global environment associations. Constructing instances from another realm can affect identity tests and \passthrough{\lstinline!instanceof!}, even if the same source-level constructor name appears (9.3, 19--20). Cross-realm equality and object brands should be reviewed when exchanging objects between contexts.

\begin{lstlisting}
// Cross-realm conceptual example; creating the second realm is HOST-specific.
// Object.prototype.toString / Array.isArray often provide better
// diagnostics than assuming a local Array constructor owns every array.
export function requireArray(value) {
  if (!Array.isArray(value)) throw new TypeError('expected array');
  return value;
}
\end{lstlisting}

\section{Completion records and propagation}\label{completion-records-and-propagation}

\textbf{{[}SPEC{]}} Algorithms describe a \passthrough{\lstinline!normal!}, \passthrough{\lstinline!throw!}, \passthrough{\lstinline!return!}, \passthrough{\lstinline!break!}, or \passthrough{\lstinline!continue!} completion where relevant; completion records carry type/value/target (5.2 and 6.2). \passthrough{\lstinline!?!} propagates abrupt completion; \passthrough{\lstinline"!"} unwraps a completion expected to be normal. In application terms, explicit \passthrough{\lstinline!return!}, thrown exception, or iterator termination have different semantics; do not flatten them into one success/failure Boolean.

\textbf{Trace template:} algorithm name; input state; each abstract operation called; user-callable hooks; normal path; abrupt path; final observable outcome. Include all conditional steps even if the chosen example takes only one branch.

\section{Ordinary and exotic object methods}\label{ordinary-and-exotic-object-methods}

\textbf{{[}SPEC{]}} Objects support essential internal methods such as \passthrough{\lstinline![[Get]]!}, \passthrough{\lstinline![[Set]]!}, \passthrough{\lstinline![[GetOwnProperty]]!}, \passthrough{\lstinline![[DefineOwnProperty]]!}, \passthrough{\lstinline![[OwnPropertyKeys]]!}, and \passthrough{\lstinline![[Call]]!} for callable objects (10). Ordinary methods implement default behavior; exotic objects, including proxies and arrays, modify rules. Proxy traps are constrained by invariants.

\begin{lstlisting}
const observed = [];
const proxy = new Proxy({ x: 1 }, {
  get(target, prop, receiver) {
    observed.push(String(prop));
    return Reflect.get(target, prop, receiver);
  }
});
console.assert(proxy.x === 1);
console.assert(observed.join(',') === 'x');
\end{lstlisting}

\textbf{{[}MODEL{]}} Accessor getters and proxy traps are potential reentrancy points. A supposedly read-only validation step can invoke user-defined code; capture snapshots before sensitive decisions if consistency is required.

\section{Script and module instantiation}\label{script-and-module-instantiation}

\textbf{{[}SPEC{]}} Scripts and source text modules have distinct parse goals, records, and instantiation/evaluation procedures (16). Module imports are live bindings. A module graph must be resolved and linked; module evaluation may be asynchronous owing to top-level \passthrough{\lstinline!await!}. Dynamic import can return a promise whose fulfillment depends on host loading plus language evaluation.

\begin{longtable}[]{@{}
  >{\raggedright\arraybackslash}p{(\columnwidth - 4\tabcolsep) * \real{0.3333}}
  >{\raggedright\arraybackslash}p{(\columnwidth - 4\tabcolsep) * \real{0.3333}}
  >{\raggedright\arraybackslash}p{(\columnwidth - 4\tabcolsep) * \real{0.3333}}@{}}
\toprule\noalign{}
\begin{minipage}[b]{\linewidth}\raggedright
Concern
\end{minipage} & \begin{minipage}[b]{\linewidth}\raggedright
ECMA-262 part
\end{minipage} & \begin{minipage}[b]{\linewidth}\raggedright
Host responsibility
\end{minipage} \\
\midrule\noalign{}
\endhead
\bottomrule\noalign{}
\endlastfoot
Grammar / module statements & 16.2 & none for valid syntax \\
Imports and module requests & 16.2 & map specifiers to modules \\
Import attributes / JSON modules & 16.2 & supported keys / resource retrieval \\
\passthrough{\lstinline!import.meta!} fields & 16.2 and Annex D & populate host-specific metadata \\
Executable module graph & 16.2 & load and supply requested modules \\
Top-level \passthrough{\lstinline!await!} & 16.2 & integrate resulting jobs into host schedule \\
\end{longtable}

\section{Promise jobs and asynchronous execution}\label{promise-jobs-and-asynchronous-execution}

\textbf{{[}SPEC{]}} Promise objects have pending/fulfilled/rejected state, reaction records, resolution procedures, and jobs for reaction execution (27.2). \passthrough{\lstinline!.then!} reactions are not called inline during \passthrough{\lstinline!.then!}; a pending fulfillment schedules work through host hooks described in 9.5. Promise states are not equivalent to whether the underlying host operation has been canceled.

\begin{lstlisting}
const order = [];
const promise = Promise.resolve('ready');
promise.then(() => order.push('then'));
order.push('sync');
await promise;
console.assert(order[0] === 'sync');
console.assert(order[1] === 'then');
\end{lstlisting}

\textbf{{[}HOST{]}} Browser microtask checkpoints, timers, input events, and Node-specific phases are host policy. Avoid claiming ECMA-262 specifies the full ordering between timer callbacks and I/O callbacks. \textbf{{[}MODEL{]}} Distinguish promise reaction jobs from event-loop phases; a diagram of one environment is not a universal language guarantee.

\section{Generators and async iteration}\label{generators-and-async-iteration}

\textbf{{[}SPEC{]}} Generators hold suspended execution states; \passthrough{\lstinline!yield!}, \passthrough{\lstinline!.next!}, \passthrough{\lstinline!.throw!}, and \passthrough{\lstinline!.return!} control reentry (15, 27). \passthrough{\lstinline!for...of!} and \passthrough{\lstinline!for await...of!} interact with iterator close semantics. ECMAScript 2025 Iterator helper objects compose iterator operations; downstream early termination can trigger upstream closure.

\begin{lstlisting}
function* ids() {
  try { yield 1; yield 2; yield 3; }
  finally { console.log('iterator cleanup'); }
}
for (const id of ids()) {
  if (id === 1) break;  // executes generator's finally
}
\end{lstlisting}

\textbf{{[}MODEL{]}} Diagram an iterator lifecycle with Created, Suspended, Executing, and Completed states; specify what happens when iteration throws or breaks. Avoid assuming all iterator adaptors behave like generators without inspection.

\section{Agents, clusters, and shared memory}\label{agents-clusters-and-shared-memory}

\textbf{{[}SPEC{]}} Agents and agent clusters appear in 9.6--9.8; shared memory and \passthrough{\lstinline!Atomics!} are defined in 25 and memory-model constraints in 29. A \emph{data race} concept is relevant to shared typed-array memory; not every concurrent access is automatically forbidden, but guarantees differ based on atomicity and ordering. Shared-memory programming should use explicit synchronization protocols.

\begin{lstlisting}
// Host is responsible for sharing buffer with a second agent.
const shared = new SharedArrayBuffer(Int32Array.BYTES_PER_ELEMENT);
const flags = new Int32Array(shared);
Atomics.store(flags, 0, 1);
console.assert(Atomics.load(flags, 0) === 1);
\end{lstlisting}

\textbf{{[}HOST{]}} Workers, shared-memory capability policy, posting messages, and OS thread realization are not fully defined by ECMA-262. \textbf{{[}MODEL{]}} For production shared-memory systems, state an invariant, a synchronization design, and a happens-before argument; test on multiple architectures.

\section{Garbage collection, WeakRef, and finalization}\label{garbage-collection-weakref-and-finalization}

\textbf{{[}SPEC{]}} \passthrough{\lstinline!WeakRef!} and \passthrough{\lstinline!FinalizationRegistry!} are for optional observation of reachability; their callbacks or dereferences cannot be used to demand deterministic cleanup. Clause 9.9 describes processing constraints, and clause 26 specifies the interfaces. \textbf{{[}MODEL{]}} Use explicit disposal or host-resource close methods for critical file handles, locks, or sensitive buffers, not finalizers.

\section{Runtime incident-analysis playbook}\label{runtime-incident-analysis-playbook}

\begin{enumerate}
\def\labelenumi{\arabic{enumi}.}
\tightlist
\item
  \textbf{Reproduce:} record engine build, host, code source type (script/module), and flags.
\item
  \textbf{Classify:} parser/early error, reference resolution, object property operation, module linkage, promise job, or host I/O.
\item
  \textbf{Trace:} map each observable step to an abstract operation or named host hook.
\item
  \textbf{Probe:} create discriminating unit tests for alternate execution paths and error completions.
\item
  \textbf{Contain:} isolate changing global objects/prototypes, plugin hooks, and shared memory.
\item
  \textbf{Conclude:} separate a language-conformance issue from implementation bug, host policy, or application defect.
\end{enumerate}

\section{Review questions and architecture deliverables}\label{review-questions-and-architecture-deliverables}

Explain the distinctions among environment, execution context, realm, job, agent, and host. Trace an async function that throws synchronously before its first \passthrough{\lstinline!await!}; describe the returned promise's rejection. Sketch module evaluation with a circular import and identify live bindings. Demonstrate a proxy trap violating an internal-method invariant. Specify atomic and nonatomic access paths for a shared flag.

\textbf{Suggested deliverables:} execution-sequence diagram, module dependency graph with async edges, binding lifetime diagram, resource-closure test case, and race-analysis checklist.

\section{Source basis}\label{source-basis}

ECMA-262:2025, 16th edition, sections 5--10, 15--16, 25--29, Annex D (uploaded original). Concepts marked \textbf{{[}MODEL{]}} are explanatory and not claimed as required implementation architecture.

\end{document}
